\chapter{Regression Discontinuity}\label{ch:b3:regression-discontinuity}

A candidate who wins an election by a hundred votes and one who loses by a
hundred are, for most purposes, the same: the same kind of district, the same
campaign, the same party. One of them takes office and the other does not. If
the party that barely won does better at the next election than the party that
barely lost, the difference is the effect of holding the seat, and almost
nothing else. A rule that switches a treatment on at a threshold --- a vote
margin of zero, a pass mark, an age, an enrolment that forces a school to open a
new class --- creates, around the threshold, something close to a randomised
experiment. This chapter shows how to read it, how to estimate the jump, and
how a design that looks this clean can still be fooled.

\begin{openingchart}{Simulated close elections: the party's vote share at the
next election against its winning margin at this one (6{,}000 elections; dots
are means over bins of 0.02 in the margin; lines are local linear fits within
0.25 of the threshold). The simulation is calibrated on a published estimate:
a jump of 0.077 in the next vote share for a party that barely
won.}{simulated data calibrated on Lee (2008), Table II; the chapter's seeded
script}
\begin{tikzpicture}
\begin{axis}[omscatter, width=0.78\linewidth, height=0.45\linewidth, xmin=-0.5, xmax=0.5,
  ymin=0, ymax=1, xlabel={Winning margin at this election},
  ylabel={Vote share at the next election},
  x tick label style={/pgf/number format/.cd, fixed, precision=2},
  y tick label style={/pgf/number format/.cd, fixed, precision=1}]
  \addplot table[x=margin, y=share] {figdata/bachelor-3/04-regression-discontinuity/out/elections_bins.csv};
  \addplot[omThm, very thick, no marks, sharp plot] table[x=margin, y=fit] {figdata/bachelor-3/04-regression-discontinuity/out/elections_fit_left.csv};
  \addplot[omThm, very thick, no marks, sharp plot] table[x=margin, y=fit] {figdata/bachelor-3/04-regression-discontinuity/out/elections_fit_right.csv};
  \draw[guide] (axis cs:0,0) -- (axis cs:0,1);
\end{axis}
\end{tikzpicture}
\end{openingchart}

The dots rise smoothly with the margin on both sides: parties that win big at
one election tend to do well at the next. At zero the curve breaks. The two
fitted lines meet the threshold at 0.4597 and 0.5353, a jump of 0.0756 in the
vote share. % figdata: elections_fit_left, elections_fit_right
The published estimate it is calibrated on is 0.077, with a standard error of
0.011, from close elections to the United States House of Representatives; the
same study finds that the party's vote share at the previous election does
\emph{not} jump at the threshold, an estimate of $-0.002$ (0.011). % ledger: b3rd:lee08.vote
That second number is the reason to believe the first: just above and just
below the threshold, the districts were alike before the election.

\begin{recall}
\Cref{ch:b3:instrumental-variables} defined the \omterm{def:b3:instrumental-variables:wald}{Wald estimator}, the ratio of
the effect of an instrument on the outcome to its effect on the treatment, and
showed that it estimates a \omterm{def:b3:instrumental-variables:late}{local average treatment effect}: the effect on the
units whose treatment the instrument changes.
\end{recall}

\begin{notation}[Symbols of this chapter]\label{not:b3:regression-discontinuity:symbols}
The \omterm{def:b3:regression-discontinuity:running-variable}{running variable} is $x$, the \omterm{def:b3:regression-discontinuity:running-variable}{cutoff} $c$, and $D_i = 1$ marks treatment.
Potential outcomes are $Y_i(1)$ and $Y_i(0)$, and $\mu_d(x) = \E[Y_i(d) \mid
x_i = x]$ their mean at a value of the \omterm{def:b3:regression-discontinuity:running-variable}{running variable}. The window is $h$.
\end{notation}

\section{A rule that jumps}\label{sec:b3:regression-discontinuity:rule}

\begin{definition}[Running variable; cutoff]\label{def:b3:regression-discontinuity:running-variable}
\index{running variable}\index{cutoff}
In a design where a rule assigns treatment according to whether an observed
variable passes a threshold, the variable is the \emph{running variable} and the
threshold the \emph{cutoff}.
\end{definition}

The examples are many because institutions love thresholds. Scholarships go to
those who score above a mark, benefits to those below an income, drinking to
those above an age, an extra class to schools above an enrolment. In each case
the units on the two sides of the \omterm{def:b3:regression-discontinuity:running-variable}{cutoff} differ in their \omterm{def:b3:regression-discontinuity:running-variable}{running variable}, and
in everything correlated with it, but only gradually; they differ in the
treatment abruptly. If nothing else jumps at the \omterm{def:b3:regression-discontinuity:running-variable}{cutoff}, a jump in the outcome
there is the treatment's effect.

\section{The sharp design}\label{sec:b3:regression-discontinuity:sharp}

\begin{definition}[Sharp regression discontinuity]\label{def:b3:regression-discontinuity:sharp}
\index{sharp regression discontinuity}
A \emph{sharp regression discontinuity} design is one in which treatment is a
deterministic function of the \omterm{def:b3:regression-discontinuity:running-variable}{running variable}: $D_i = 1$ if $x_i \ge c$ and
$D_i = 0$ otherwise. Its estimand is the jump $\lim_{x \downarrow c} \E[Y_i
\mid x_i = x] - \lim_{x \uparrow c} \E[Y_i \mid x_i = x]$.
\end{definition}

\begin{proposition}[Continuity identifies the effect at the cutoff]\label{prop:b3:regression-discontinuity:continuity}
In a sharp design, if $\mu_0(x)$ and $\mu_1(x)$ are continuous at $x = c$, the
jump in the mean outcome at the \omterm{def:b3:regression-discontinuity:running-variable}{cutoff} equals the average treatment effect at
the \omterm{def:b3:regression-discontinuity:running-variable}{cutoff}, $\mu_1(c) - \mu_0(c) = \E[Y_i(1) - Y_i(0) \mid x_i = c]$.
\end{proposition}

\begin{proof}
For $x \ge c$ every unit is treated, so $\E[Y_i \mid x_i = x] = \mu_1(x)$; for $x <
c$, $\E[Y_i \mid x_i = x] = \mu_0(x)$. The limit from above is $\mu_1(c)$ and
the limit from below $\mu_0(c)$, both by continuity. Their difference is
$\mu_1(c) - \mu_0(c)$, which is $\E[Y_i(1) - Y_i(0) \mid x_i = c]$ by the
definition of $\mu_d$.
\end{proof}

The proof uses nothing about how units came to be on either side: they can
have chosen their \omterm{def:b3:regression-discontinuity:running-variable}{running variable}, as candidates choose how hard to
campaign. It needs only that they could not choose it \emph{precisely}. A
candidate can raise the expected margin, but the count on the night has a
random element, so that among races within a few votes of a tie, which side
won is as good as a coin toss. When units can sort precisely around the \omterm{def:b3:regression-discontinuity:running-variable}{cutoff},
the continuity of $\mu_0$ and $\mu_1$ fails, and so does the design
(\cref{sec:b3:regression-discontinuity:manipulation}).

The other price is in the estimand: the effect at the \omterm{def:b3:regression-discontinuity:running-variable}{cutoff}, for the units
there. The effect of incumbency on parties that win by forty points, or of a
scholarship on pupils who score twenty marks above the line, is not
identified.

\section{Estimation: local linear fits and the window}\label{sec:b3:regression-discontinuity:estimation}

The jump is a difference between two limits, and data are never at the limit.
The estimator must use observations some distance from the \omterm{def:b3:regression-discontinuity:running-variable}{cutoff}, and it must
not mistake the slope of the curve for a jump.

\begin{definition}[Window; local linear regression]\label{def:b3:regression-discontinuity:local-linear}
\index{window (regression discontinuity)}\index{local linear regression}
The \emph{window} (or bandwidth) $h$ is the distance from the \omterm{def:b3:regression-discontinuity:running-variable}{cutoff} within which
observations are used. A \emph{local linear regression} estimates each limit by
the intercept, at the \omterm{def:b3:regression-discontinuity:running-variable}{cutoff}, of a weighted least-squares line fitted on that
side of the \omterm{def:b3:regression-discontinuity:running-variable}{cutoff} to the observations within the window. With triangular
weights, an observation at distance $|x_i - c|$ gets weight $1 - |x_i -
c|/h$.
\end{definition}

The simplest estimator, a \emph{local constant}, compares the mean outcomes on
the two sides within the window. If the curve slopes upward, the mean just above
the \omterm{def:b3:regression-discontinuity:running-variable}{cutoff} sits above the curve's value at the \omterm{def:b3:regression-discontinuity:running-variable}{cutoff}, and the mean just below
sits beneath it: the difference of means overstates the jump by about the slope
times the window. A local linear fit removes that first-order error, leaving one
that depends on the curvature, which is of the order of $h^2$.

\admitted
The precise statements --- the bias of the local constant estimator is
proportional to $h$ and that of the local linear estimator to $h^2$, under
smoothness of $\mu_0$ and $\mu_1$ --- are Taylor expansions of the two limits;
the variance of both is proportional to $1/(nh)$, the number of observations
in the window.

\omcode{code/bachelor-3/04-regression-discontinuity/listings/rd_sharp.py}{1}{18}{A sharp design on the simulated elections; \texttt{local\_linear} fits each side with triangular weights.}\label{lst:b3:regression-discontinuity:sharp}

\omoutput{code/bachelor-3/04-regression-discontinuity/outputs/rd_sharp.txt}{True jump 0.077; robust standard errors.}\label{out:b3:regression-discontinuity:sharp}

The three local linear estimates are 0.078, 0.081 and 0.076; the difference of
means within 0.05 is 0.102, a third too large. % output: rd_sharp
The last line is the balance test: the vote share at the \emph{previous}
election, decided before anything happened at the threshold, shows no jump.

\begin{definition}[Manipulation test; covariate balance test]\label{def:b3:regression-discontinuity:tests}
\index{covariate balance test}\index{manipulation test}
A \emph{covariate balance test} estimates the jump at the \omterm{def:b3:regression-discontinuity:running-variable}{cutoff} in a variable
determined before treatment; under the design's assumptions it is zero. A
\emph{manipulation test} (or density test) estimates the jump at the \omterm{def:b3:regression-discontinuity:running-variable}{cutoff} in
the density of the \omterm{def:b3:regression-discontinuity:running-variable}{running variable}; a design in which units cannot sort
precisely has none.
\end{definition}

How wide a window? A wide one uses more data and is less noisy; a narrow one is
less biased. The chart shows the trade-off on 300 simulated samples of the
opening chart's elections.

\begin{omchart}{Estimates of a jump of 0.077 against the window, on 300 simulated
samples of 6{,}000 elections: mean estimate (lines) and one standard deviation
either side (bands), for the difference of means and the local linear
fit.}{simulated data, the chapter's seeded script}
\begin{tikzpicture}
\begin{axis}[omtime, width=0.78\linewidth, height=0.42\linewidth, xmin=0, xmax=0.52,
  ymin=0.04, ymax=0.21, xlabel={Window $h$ (vote margin)}, ylabel={Estimated jump},
  legend pos=north west, xtick={0,0.1,0.2,0.3,0.4,0.5},
  x tick label style={/pgf/number format/.cd, fixed, precision=1},
  y tick label style={/pgf/number format/.cd, fixed, precision=2}]
  \addplot[name path=cu, draw=none, forget plot] table[x=h, y expr=\thisrow{const_mean}+\thisrow{const_sd}] {figdata/bachelor-3/04-regression-discontinuity/out/bandwidth.csv};
  \addplot[name path=cl, draw=none, forget plot] table[x=h, y expr=\thisrow{const_mean}-\thisrow{const_sd}] {figdata/bachelor-3/04-regression-discontinuity/out/bandwidth.csv};
  \addplot[omMeth!25, forget plot] fill between[of=cu and cl];
  \addplot[name path=lu, draw=none, forget plot] table[x=h, y expr=\thisrow{lin_mean}+\thisrow{lin_sd}] {figdata/bachelor-3/04-regression-discontinuity/out/bandwidth.csv};
  \addplot[name path=ll, draw=none, forget plot] table[x=h, y expr=\thisrow{lin_mean}-\thisrow{lin_sd}] {figdata/bachelor-3/04-regression-discontinuity/out/bandwidth.csv};
  \addplot[omDef!25, forget plot] fill between[of=lu and ll];
  \addplot[omMeth, very thick, no marks] table[x=h, y=const_mean] {figdata/bachelor-3/04-regression-discontinuity/out/bandwidth.csv};
  \addplot[omDef, very thick, no marks] table[x=h, y=lin_mean] {figdata/bachelor-3/04-regression-discontinuity/out/bandwidth.csv};
  \draw[guide] (axis cs:0,0.077) -- (axis cs:0.52,0.077);
  \legend{difference of means, local linear}
\end{axis}
\end{tikzpicture}
\end{omchart}

The difference of means drifts away from the truth almost linearly as the
window widens; the local linear fit stays on it until about 0.2 and then
bends down as the curvature of the cubic catches it. Its band, wide at small
windows, narrows steadily. \Cref{lst:b3:regression-discontinuity:bandwidth}
puts numbers on the picture.

\omcode{code/bachelor-3/04-regression-discontinuity/listings/rd_bandwidth.py}{1}{21}{Bias, noise and their combination against the window.}\label{lst:b3:regression-discontinuity:bandwidth}

\omoutput{code/bachelor-3/04-regression-discontinuity/outputs/rd_bandwidth.txt}{500 samples; true jump 0.077.}\label{out:b3:regression-discontinuity:bandwidth}

The local linear fit has its smallest root mean squared error, 0.006, at a
window of 0.25; the difference of means has its own at the narrowest window,
0.02, and even there it does worse than the local linear fit at its best.
% output: rd_bandwidth

\begin{measure}[Choosing the window]
The root mean squared error trades a bias that grows with $h$ against a variance
that shrinks with it, and its minimum depends on the curvature of $\mu_0$ and
$\mu_1$ near the \omterm{def:b3:regression-discontinuity:running-variable}{cutoff} and on the noise --- both estimable. Imbens and
Kalyanaraman (2012) derived a window that minimises an estimate of that error;
Calonico, Cattaneo and Titiunik (2014) showed that intervals built around an
error-minimising window are centred on a biased estimate and cover the truth
too rarely, and corrected them by estimating and subtracting the bias, with a
standard error that accounts for the correction. Two habits follow. Report the
estimate at the chosen window and at half and twice that window: an estimate
that moves a great deal is not one to trust. And never fit one high-order
polynomial to all the data on each side; such fits give large weight to
observations far from the \omterm{def:b3:regression-discontinuity:running-variable}{cutoff} and produce erratic jumps
(\cref{exo:b3:regression-discontinuity:5}).
\end{measure}

\section{The fuzzy design}\label{sec:b3:regression-discontinuity:fuzzy}

Many rules do not bind exactly. A scholarship is offered above a mark, but not
everyone takes it up, and some below the mark find funding elsewhere. Treatment
then jumps at the \omterm{def:b3:regression-discontinuity:running-variable}{cutoff} by less than one.

\begin{definition}[Fuzzy regression discontinuity]\label{def:b3:regression-discontinuity:fuzzy}
\index{fuzzy regression discontinuity}
A \emph{fuzzy regression discontinuity} design is one in which the probability of
treatment jumps at the \omterm{def:b3:regression-discontinuity:running-variable}{cutoff} by less than one: $0 < \lim_{x \downarrow c}
\Pr(D_i = 1 \mid x) - \lim_{x \uparrow c} \Pr(D_i = 1 \mid x) < 1$. Its estimand
is the ratio of the jump in the mean outcome to the jump in the probability of
treatment.
\end{definition}

\begin{omfigure}
\begin{tikzpicture}[econ, scale=0.8]
  \begin{scope}
    \draw[econ axes] (0,3.2) node[above] {$\Pr(D = 1)$} |- (5.2,0) node[right] {$x$};
    \draw[guide] (2.6,0) node[below] {$c$} -- (2.6,3);
    \draw[curve, omDef] (0,0.05) -- (2.6,0.05);
    \draw[curve, omDef] (2.6,2.6) -- (5,2.6);
    \node[font=\footnotesize] at (2.5,-0.9) {sharp};
    \node[font=\scriptsize, left] at (0,2.6) {1};
  \end{scope}
  \begin{scope}[xshift=7cm]
    \draw[econ axes] (0,3.2) node[above] {$\Pr(D = 1)$} |- (5.2,0) node[right] {$x$};
    \draw[guide] (2.6,0) node[below] {$c$} -- (2.6,3);
    \draw[curve, omThm] (0,0.1) .. controls (1.5,0.2) .. (2.6,0.35);
    \draw[curve, omThm] (2.6,1.9) .. controls (3.8,2.1) .. (5,2.25);
    \draw[<->, thick] (2.85,0.35) -- node[right, font=\scriptsize] {jump $<1$} (2.85,1.9);
    \node[font=\footnotesize] at (2.5,-0.9) {fuzzy};
    \node[font=\scriptsize, left] at (0,2.6) {1};
    \draw[guide] (0,2.6) -- (5,2.6);
  \end{scope}
\end{tikzpicture}
\par\smallskip{\small The probability of treatment against the \omterm{def:b3:regression-discontinuity:running-variable}{running variable}.
In a sharp design it jumps from 0 to 1 at the \omterm{def:b3:regression-discontinuity:running-variable}{cutoff}; in a fuzzy design it jumps
by less.}
\end{omfigure}

\begin{proposition}[Fuzzy regression discontinuity as a Wald ratio]\label{prop:b3:regression-discontinuity:fuzzy}
Suppose that near the \omterm{def:b3:regression-discontinuity:running-variable}{cutoff} the jump in treatment is caused by the rule, that
crossing the \omterm{def:b3:regression-discontinuity:running-variable}{cutoff} affects the outcome only through treatment, that no unit is
treated below the \omterm{def:b3:regression-discontinuity:running-variable}{cutoff} and untreated above it who would have been the
opposite way round (monotonicity), and that the mean potential outcomes and
treatment shares of each type of unit are continuous at $c$. Then the fuzzy
estimand equals the average treatment effect at the \omterm{def:b3:regression-discontinuity:running-variable}{cutoff} for the \omterm{def:b3:instrumental-variables:compliance}{compliers},
the units treated if and only if they are above it.
\end{proposition}

\begin{proof}
Let $Z_i = 1$ if $x_i \ge c$: $Z$ is an instrument for $D$ for the units at the
\omterm{def:b3:regression-discontinuity:running-variable}{cutoff}. Taking limits from each side plays the role of comparing $Z = 1$ with
$Z = 0$ among units with the same $x = c$. The proof of the LATE theorem in
\cref{ch:b3:instrumental-variables} then applies term by term with
expectations replaced by these limits: the jump in $\E[Y]$ is the share of
\omterm{def:b3:instrumental-variables:compliance}{compliers} at $c$ times their mean effect, the jump in $\Pr(D = 1)$ is the share
of \omterm{def:b3:instrumental-variables:compliance}{compliers}, and the ratio is their mean effect.
\end{proof}

The classic fuzzy design is in schools. Angrist and Lavy (1999) noticed that
Israeli schools split a grade into classes of at most 40, so that a grade of 40
pupils sits in one class of 40 and a grade of 41 in two classes of 20 or 21.
% ledger: b3rd:al.rule
Class size predicted by the rule falls abruptly at enrolments of 41, 81 and
121, and actual class size follows, though not exactly.

\begin{omchart}{Class size against grade enrolment. Line: the size predicted by a
cap of 40 per class, $r/(\lfloor (r-1)/40 \rfloor + 1)$. Dots: mean class size at
each enrolment in 4{,}000 simulated schools, three quarters of which follow the
rule and the rest a target size of their own.}{the rule as stated in
Angrist, Lavy, Leder-Luis and Shany (2019); simulated schools, the chapter's
seeded script}
\begin{tikzpicture}
\begin{axis}[omtime, width=0.78\linewidth, height=0.4\linewidth, xmin=0, xmax=162,
  ymin=0, ymax=45, xlabel={Pupils enrolled in the grade}, ylabel={Pupils per class},
  xtick={0,40,80,120,160}, legend pos=south east]
  \addplot[omDef, thick, no marks] table[x=enrol, y=rule] {figdata/bachelor-3/04-regression-discontinuity/out/rule.csv};
  \addplot[omThm, only marks, mark=*, mark size=1.1pt] table[x=enrol, y=size] {figdata/bachelor-3/04-regression-discontinuity/out/schools.csv};
  \legend{the rule, simulated schools}
\end{axis}
\end{tikzpicture}
\end{omchart}

\omcode{code/bachelor-3/04-regression-discontinuity/listings/rd_fuzzy.py}{1}{22}{A fuzzy design at the first \omterm{def:b3:regression-discontinuity:running-variable}{cutoff}: \omterm{def:b3:instrumental-variables:tsls}{first stage}, \omterm{def:b3:instrumental-variables:tsls}{reduced form} and their ratio.}\label{lst:b3:regression-discontinuity:fuzzy}

\omoutput{code/bachelor-3/04-regression-discontinuity/outputs/rd_fuzzy.txt}{True effect $-0.25$ points per pupil.}\label{out:b3:regression-discontinuity:fuzzy}

Crossing from 40 to 41 pupils cuts class size by 14.0 pupils and raises the mean
score by 4.0 points; the ratio is $-0.29$ points per extra pupil, against a true
$-0.25$. % output: rd_fuzzy
(From the printed numbers, $4.04/(-14.01) = -0.288$; the script divides the
unrounded jumps.)

\begin{record}[Class size in Israel]
In 1991 Israeli test data, Angrist and Lavy (1999) found that the up-and-down
pattern of class size induced by the rule was matched by a down-and-up pattern
in test scores: smaller classes raised the reading and mathematics scores of
fifth graders, with smaller gains for fourth graders and little for third
graders. % ledger: b3rd:al99.result
Twenty years later the same authors and two colleagues repeated the study on
Israeli data from 2002--2011. The rule still moved class size strongly, but
the estimated effect of class size on achievement was now zero. They also found
too many schools with enrolments just above the \omterm{def:b3:regression-discontinuity:running-variable}{cutoffs} --- one pupil more
buys a whole extra class --- and checked that this sorting did not explain
either the old effect or the new zero. % ledger: b3rd:alls.redux
\end{record}

\section{Can the running variable be manipulated?}\label{sec:b3:regression-discontinuity:manipulation}

The design's weak point is sorting. If those just above the \omterm{def:b3:regression-discontinuity:running-variable}{cutoff} got there on
purpose --- a pupil who pleads for one more mark, a school that enrols one more
child, a household that earns a little less to stay under a benefit threshold
--- then those just above differ from those just below in their wish and
ability to sort, and continuity fails. Sorting leaves a trace: too many units
just above the \omterm{def:b3:regression-discontinuity:running-variable}{cutoff}, too few just below.

\begin{omchart}{Exam scores (whole marks) of 20{,}000 simulated candidates, with
a pass mark of 70, after a marker has moved half of the scripts on 66--69 up to
70. Bars: the re-marked scores; line: the scores before re-marking.}{simulated
data, the chapter's seeded script}
\begin{tikzpicture}
\begin{axis}[omtime, width=0.78\linewidth, height=0.38\linewidth, xmin=29, xmax=96,
  ymin=0, ymax=1700, xlabel={Score}, ylabel={Candidates}, legend pos=north east,
  xtick={30,40,...,90}]
  \addplot[ybar, area legend, bar width=3.5pt, fill=omDef!45, draw=omDef] table[x=score, y=remarked] {figdata/bachelor-3/04-regression-discontinuity/out/scores.csv};
  \addplot[omThm, thick, sharp plot, no marks] table[x=score, y=clean] {figdata/bachelor-3/04-regression-discontinuity/out/scores.csv};
  \draw[guide] (axis cs:69.5,0) -- (axis cs:69.5,1700);
  \legend{re-marked, before}
\end{axis}
\end{tikzpicture}
\end{omchart}

The density test of McCrary (2008) estimates the density of the \omterm{def:b3:regression-discontinuity:running-variable}{running variable}
just above and just below the \omterm{def:b3:regression-discontinuity:running-variable}{cutoff} by local linear fits to a histogram, and
tests whether the log of their ratio is zero.

\omcode{code/bachelor-3/04-regression-discontinuity/listings/density_test.py}{1}{15}{The density test on the clean and the re-marked scores.}\label{lst:b3:regression-discontinuity:density}

\omoutput{code/bachelor-3/04-regression-discontinuity/outputs/density_test.txt}{Log ratio of the densities just above and just below 70.}\label{out:b3:regression-discontinuity:density}

On the clean scores the test finds nothing, $t = -1.1$; after 982 scripts are
moved, the log density jumps by 1.86, $t = 29.9$. % output: density_test
Real manipulation is rarely so blatant, and the test has an obvious blind
spot: if some units are pushed up and an equal number pushed down, the density
is smooth while the composition on each side is not. Balance tests on
predetermined covariates complement it. Neither can prove the design valid; a
failure of either is hard to explain away. When sorting is confined to a few
units at the \omterm{def:b3:regression-discontinuity:running-variable}{cutoff}, a \emph{donut} design drops the observations closest to it
and estimates from the rest (\cref{exo:b3:regression-discontinuity:8}).

\section{Thresholds in practice}\label{sec:b3:regression-discontinuity:applications}

\begin{record}[The minimum legal drinking age]
In the United States the minimum legal drinking age is 21, and age is a \omterm{def:b3:regression-discontinuity:running-variable}{running
variable} nobody can manipulate. Carpenter and Dobkin (2009) found that at 21 the
number of days on which young adults drink jumps by 21~\%, and the mortality
rate jumps by 9~\%. % ledger: b3rd:cd09.mortality
The design answers one question cleanly: what happens to people who turn 21
under a law of 21. Whether a different age would change total mortality, or
only move it, is a question the threshold cannot answer.
\end{record}

\begin{model}{Regression discontinuity}
\assumes Treatment, or its probability, jumps at a known \omterm{def:b3:regression-discontinuity:running-variable}{cutoff} of an observed
\omterm{def:b3:regression-discontinuity:running-variable}{running variable}; units cannot place themselves precisely on one side; the
mean potential outcomes are continuous at the \omterm{def:b3:regression-discontinuity:running-variable}{cutoff}; in a fuzzy design,
crossing the \omterm{def:b3:regression-discontinuity:running-variable}{cutoff} affects the outcome only through treatment.
\predicts A jump in the outcome at the \omterm{def:b3:regression-discontinuity:running-variable}{cutoff} equal to the effect there; no jump
in predetermined covariates; a smooth density of the \omterm{def:b3:regression-discontinuity:running-variable}{running variable}; no jump at
placebo \omterm{def:b3:regression-discontinuity:running-variable}{cutoffs} where nothing changes.
\tested In close elections, the vote share before the election does not jump
while the one after does; histograms of the \omterm{def:b3:regression-discontinuity:running-variable}{running variable} have caught real
sorting, as at enrolment \omterm{def:b3:regression-discontinuity:running-variable}{cutoffs} in Israeli schools.
\breaks When units sort precisely (enrolments, re-marked exams, income just under
a threshold); when another policy uses the same \omterm{def:b3:regression-discontinuity:running-variable}{cutoff} (the age of 21 may also
change other rules); when the curve is fitted too far from the \omterm{def:b3:regression-discontinuity:running-variable}{cutoff}; and
for questions about units far from the \omterm{def:b3:regression-discontinuity:running-variable}{cutoff}, where it says nothing.
\end{model}

\section{Exercises}\label{sec:b3:regression-discontinuity:exercises}

\begin{exercise}[$\star$]\label{exo:b3:regression-discontinuity:1}
Read the opening chart: the values of the two fitted lines at the \omterm{def:b3:regression-discontinuity:running-variable}{cutoff} and
the jump. Why are the dots near $\pm 0.4$ more scattered than those near the
\omterm{def:b3:regression-discontinuity:running-variable}{cutoff}?
\end{exercise}

\begin{exercise}[$\star$]\label{exo:b3:regression-discontinuity:2}
Compute the class size predicted by Maimonides' rule at enrolments of 40, 41, 80,
81 and 121. At which enrolments does it jump, and by how much at each?
\end{exercise}

\begin{exercise}[$\star$]\label{exo:b3:regression-discontinuity:3}
Among applicants scoring 68 or 69 on an entrance test, 40\,\% find a job within a
year; among those scoring 70 or 71, admitted to a training scheme, 52\,\%.
Estimate the scheme's effect on employment. What does the estimate assume, and
why might it overstate the effect if the employment rate rises with the score?
\end{exercise}

\begin{exercise}[$\star$]\label{exo:b3:regression-discontinuity:4}
A balance test finds that parents' income jumps by 8~\% at a scholarship's
\omterm{def:b3:regression-discontinuity:running-variable}{cutoff}. What does that suggest, and what would you look at next?
\end{exercise}

\begin{exercise}[$\star\star$]\label{exo:b3:regression-discontinuity:5}
On 300 samples of 800 simulated elections, compare global polynomials of order
1, 2, 4 and 6 on each side, fitted to all the data, with a local linear fit
within 0.25. Report the bias and the standard deviation of each, and say which
you would use.
\end{exercise}

\begin{exercise}[$\star\star$]\label{exo:b3:regression-discontinuity:6}
At an age threshold, the share of people with health insurance jumps from 0.60
to 0.85 and the share who visit a doctor in a year jumps from 0.70 to 0.75.
\begin{enumerate}
\item Estimate the effect of insurance on visits, and say for whom.
\item Why is the ratio not the effect of insurance on the 60~\% already insured
below the threshold?
\end{enumerate}
\end{exercise}

\begin{exercise}[$\star\star$]\label{exo:b3:regression-discontinuity:7}
Students who score just below a scholarship's \omterm{def:b3:regression-discontinuity:running-variable}{cutoff} can ask for their exam to
be re-marked, and the most determined ones do. If determination raises later
earnings, in which direction is a sharp estimate of the scholarship's effect
biased?
\end{exercise}

\begin{exercise}[$\star\star$]\label{exo:b3:regression-discontinuity:8}
Give the re-marked exams of \cref{lst:b3:regression-discontinuity:density} an
outcome: log earnings that rise with the true score by 0.01 per mark and jump by
0.10 for those who pass. Estimate the jump with all the recorded scores and with
donuts that drop the one and two marks on each side of the \omterm{def:b3:regression-discontinuity:running-variable}{cutoff}. Explain the
results.
\end{exercise}

\begin{exercise}[$\star\star$]\label{exo:b3:regression-discontinuity:9}
A city announces that households below an income line will receive a subsidy,
but sets the line only after the year's incomes have been reported. What pattern
in the data would suggest the line was chosen to include particular
households, and why would that matter for a regression discontinuity study?
\end{exercise}

\begin{exercise}[$\star\star\star$]\label{exo:b3:regression-discontinuity:10}
Show that with uniform weights the local linear estimate of the limit from
above, $\hat a_1$, is a weighted mean of the outcomes in the window above the
\omterm{def:b3:regression-discontinuity:running-variable}{cutoff}, $\hat a_1 = \sum_i \omega_i Y_i$, with weights $\omega_i = (S_2 - S_1
x_i)/(n_1 S_2 - S_1^2)$, where $x_i$ is measured from the \omterm{def:b3:regression-discontinuity:running-variable}{cutoff}, $n_1$ is the
number of observations and $S_k = \sum_i x_i^k$. Show that the weights sum to
one, that some are negative, and which.
\end{exercise}

\begin{exercise}[$\star\star\star$]\label{exo:b3:regression-discontinuity:11}
On 500 simulated samples of the opening chart's elections, compute how often the
95~\% interval of the local linear estimate covers the true jump, at windows
0.02, 0.10, 0.25 and 0.50. Explain the pattern using \cref{out:b3:regression-discontinuity:bandwidth}.
\end{exercise}

\begin{exercise}[$\star\star\star$]\label{exo:b3:regression-discontinuity:12}
A regression discontinuity study finds that a scholarship offered from a score
of 70 raises earnings by 12~\% for pupils at 70. A minister proposes lowering
the \omterm{def:b3:regression-discontinuity:running-variable}{cutoff} to 50. Give two reasons why the effect at 50 might differ, and
describe a design that could estimate it.
\end{exercise}

\section{Problem: The Brenmoor Scholarship}\label{sec:b3:regression-discontinuity:problem}

\begin{problem}[{Weekend problem --- the Brenmoor scholarship}]\label{pb:b3:regression-discontinuity:scholarship}
Brenmoor, the capital of Brenland (invented), offers a scholarship to every
school-leaver who scores 70 or more out of 100 in the final exam. Scores are
whole marks. A study follows 6{,}000 school-leavers to the age of 30 and records
their log earnings, whether they received the scholarship or an equivalent award
from elsewhere, and their parents' log income at the time of the exam. A
statistician fits \omterm{def:b3:regression-discontinuity:local-linear}{local linear regressions} with triangular weights, with the
\omterm{def:b3:regression-discontinuity:running-variable}{cutoff} placed at 69.5, first for all pupils and then for all pupils except those
of one school, where, it emerged, a teacher had re-marked scripts, moving some
scripts marked 67--69 up to 70. Robust
standard errors are in parentheses. % figdata: scholarship

\begin{center}\small
\pgfplotstabletypeset[col sep=comma, string type,
  columns={quantity,window,all,without},
  columns/quantity/.style={column name={Jump in}, column type=l},
  columns/window/.style={column name={Window}, column type=c},
  columns/all/.style={column name={All pupils}, column type=l},
  columns/without/.style={column name={Without the school}, column type=l},
  every head row/.style={after row=\hline}, every last row/.style={after row=\hline}]
  {figdata/bachelor-3/04-regression-discontinuity/out/scholarship.csv}
\end{center}
``Earnings'' and ``award'' are at the \omterm{def:b3:regression-discontinuity:running-variable}{cutoff} of 70; ``earnings at 60'' places
the \omterm{def:b3:regression-discontinuity:running-variable}{cutoff} at 59.5 instead, where nothing happens. A density test at the \omterm{def:b3:regression-discontinuity:running-variable}{cutoff}
(window 10) gives a log jump in the density of scores of 0.3036 (0.0798) for all
pupils and 0.0623 (0.0841) without the school. % figdata: scholarship_density

\medskip\noindent\textbf{Part I --- the offer.}
\begin{enumerate}
\item Is the design sharp or fuzzy for the offer? For the award?
\item Interpret the jump in earnings at a window of 10, for all pupils, as an
effect of being offered the scholarship.
\item Give its 95~\% confidence interval. Is the effect of the offer
distinguishable from zero?
\item Compare the estimates at windows 5, 10 and 20. What does their stability
suggest, and why do the standard errors fall?
\item Interpret the jump in parents' income. Why is it reassuring?
\item Why is the effect of the offer smaller than the effect of the
scholarship itself?
\end{enumerate}

\medskip\noindent\textbf{Part II --- take-up.}
\begin{enumerate}[resume]
\item Interpret the jump in the award at a window of 10.
\item Compute the fuzzy estimate of the effect of the award on log earnings at
a window of 10, and convert it into a percentage.
\item Compute it at windows 5 and 20.
\item Approximate the standard error of the fuzzy estimate at a window of 10 by
dividing the standard error of the earnings jump by the award jump, and give a
95~\% interval.
\item For whom is the fuzzy estimate an effect? Describe the two kinds of pupil
it leaves out.
\item State the assumption that the offer affects earnings only through the
award, and give one way it could fail.
\end{enumerate}

\medskip\noindent\textbf{Part III --- checks.}
\begin{enumerate}[resume]
\item Interpret the jump in earnings at 60.
\item The placebo uses scores from 50 to 69 only. Why must the window of a
placebo \omterm{def:b3:regression-discontinuity:running-variable}{cutoff} stop short of the real \omterm{def:b3:regression-discontinuity:running-variable}{cutoff}?
\item At what window would you report the main estimate, and what else would
you report beside it?
\item The study's critics say the scholarship's effect is mostly the effect of
passing a prestigious threshold, which employers see. What evidence could
separate the two?
\end{enumerate}

\medskip\noindent\textbf{Part IV --- the teacher.}
\begin{enumerate}[resume]
\item Compute the $t$-statistic of the density test for all pupils and without
the school. What do they show?
\item Convert the log jump for all pupils into a ratio of densities.
\item Who are the re-marked pupils, and how do they differ from the others just
above 70?
\item In which direction should the re-marking bias the jump in earnings? Compare
the estimates with and without the school at each window.
\item Why does the jump in parents' income not reveal the re-marking?
\item Dropping the school removes the manipulation but also some honest pupils.
When would that be a problem?
\item Propose a donut design for the full sample and say what it would drop.
\item State the result: the effect of the scholarship on earnings at the
threshold, with its interval and the evidence from the density test.
\end{enumerate}
\end{problem}
